% ============================================================
%  ElectroLab-JS ? Artículo Científico
%  Formato: IEEE / Artículo de conferencia académica
%  Motor: pdfLaTeX  |  Codificación: UTF-8
% ============================================================
\documentclass[12pt, a4paper]{article}

% ?? Paquetes base ????????????????????????????????????????????
\usepackage[utf8]{inputenc}
\usepackage[T1]{fontenc}
\usepackage[provide=*,spanish]{babel}

\usepackage[protrusion=true,expansion=false]{microtype}

% ?? Geometría y márgenes ?????????????????????????????????????
\usepackage[
  top=2.5cm, bottom=2.5cm,
  left=3.0cm, right=2.5cm,
  headheight=14pt
]{geometry}

% ?? Gráficos y color ?????????????????????????????????????????
\usepackage{graphicx}
\usepackage[dvipsnames, table]{xcolor}
\usepackage{float}
\usepackage{caption}
\usepackage{subcaption}

% ?? Tablas ???????????????????????????????????????????????????
\usepackage{booktabs}
\usepackage{tabularx}
\usepackage{array}
\usepackage{multirow}
\usepackage{longtable}

% ?? Matemáticas ??????????????????????????????????????????????
\usepackage{amsmath}
\usepackage{amssymb}
\usepackage{mathtools}
\usepackage{bm}
\usepackage{siunitx}
\sisetup{group-separator={,}}

% ?? Código fuente ????????????????????????????????????????????
\usepackage{listings}
\usepackage{lstautogobble}

% ?? Cajas de resaltado ???????????????????????????????????????
\usepackage{tcolorbox}
\tcbuselibrary{breakable, skins, most}

% ?? Encabezados, pies y secciones ????????????????????????????
\usepackage{fancyhdr}
\usepackage{titlesec}
\usepackage{titletoc}
\usepackage{enumitem}
\usepackage{parskip}

% ?? Hipervínculos y referencias ??????????????????????????????
\usepackage[
  colorlinks=true,
  linkcolor=NavyBlue,
  citecolor=NavyBlue,
  urlcolor=NavyBlue,
  pdftitle={ElectroLab-JS: Simulador Interactivo de Electricidad y Magnetismo},
  pdfauthor={[Tu Nombre Completo]},
]{hyperref}
\usepackage{cite}
\usepackage{url}

% ?? Marcadores de posición para figuras ??????????????????????
\usepackage{tocloft}
\usepackage{afterpage}

% ============================================================
%  COLORES PERSONALIZADOS
% ============================================================
\definecolor{azulprimario}{HTML}{1F3864}
\definecolor{azulsecundario}{HTML}{2E75B6}
\definecolor{verdecaja}{HTML}{1F4E39}
\definecolor{verdebg}{HTML}{E2EFDA}
\definecolor{verdeborde}{HTML}{70AD47}
\definecolor{amarillocaja}{HTML}{5C4000}
\definecolor{amarillobg}{HTML}{FFF2CC}
\definecolor{amarilloborde}{HTML}{C9A800}
\definecolor{griscode}{HTML}{F5F5F5}
\definecolor{azulcod}{HTML}{1F3864}
\definecolor{grisclaro}{HTML}{F0F4F8}

% ============================================================
%  CONFIGURACIÓN DE SECCIONES
% ============================================================
\titleformat{\section}
  {\Large\bfseries\color{azulprimario}}
  {\thesection.}{0.6em}{}
  [\vspace{-0.3em}\textcolor{azulsecundario}{\rule{\linewidth}{1.5pt}}]

\titleformat{\subsection}
  {\large\bfseries\color{azulsecundario}}
  {\thesubsection}{0.6em}{}

\titleformat{\subsubsection}
  {\normalsize\bfseries\color{black!70}}
  {\thesubsubsection}{0.6em}{}

\titlespacing*{\section}{0pt}{18pt}{8pt}
\titlespacing*{\subsection}{0pt}{12pt}{5pt}
\titlespacing*{\subsubsection}{0pt}{8pt}{3pt}

% ============================================================
%  CAJAS DE ENTORNO
% ============================================================

% Caja "Experiencia de clase"
\newtcolorbox{cajaclas}[1][]{
  enhanced, breakable,
  colback=verdebg, colframe=verdeborde,
  left=6pt, right=6pt, top=5pt, bottom=5pt,
  fontupper=\small\itshape\color{verdecaja},
  title={\small\textbf{Experiencia de laboratorio / clase}},
  attach boxed title to top left={yshift=-2mm, xshift=4mm},
  boxed title style={colback=verdeborde, colframe=verdeborde},
  #1
}

% Caja "Nota histórica / dato curioso"
\newtcolorbox{cajahist}[1][]{
  enhanced, breakable,
  colback=amarillobg, colframe=amarilloborde,
  left=6pt, right=6pt, top=5pt, bottom=5pt,
  fontupper=\small\itshape\color{amarillocaja},
  borderline west={3pt}{0pt}{amarilloborde},
  #1
}

% Caja de caso de uso
\newtcolorbox{casodeuso}[2][]{
  enhanced, breakable,
  colback=grisclaro, colframe=azulsecundario,
  left=8pt, right=8pt, top=6pt, bottom=6pt,
  title={\textbf{ Caso de uso: #2}},
  fonttitle=\normalsize\bfseries\color{white},
  attach boxed title to top left={yshift=-2mm, xshift=4mm},
  boxed title style={colback=azulsecundario},
  #1
}

% ============================================================
%  PLACEHOLDER DE FIGURA
% ============================================================
\newcommand{\figplaceholder}[2]{%
  \begin{figure}[H]
    \centering
    \begin{tcolorbox}[
      colback=white, colframe=gray!40,
      width=0.9\linewidth,
      halign=center, valign=center,
      height=4.5cm,
      enhanced, sharp corners,
      borderline={0.5pt}{0pt}{gray!40, dashed}
    ]
      \centering
      \textcolor{gray!60}{\Large } \\[4pt]
      \textbf{\textcolor{gray!70}{#1}} \\[2pt]
      \textit{\small\textcolor{gray!60}{#2}}
    \end{tcolorbox}
    \caption{#1}
    \label{fig:#1}
  \end{figure}%
}

% ============================================================
%  LISTINGS (código JS)
% ============================================================
\lstdefinelanguage{JavaScript}{
  keywords={const,let,var,function,return,if,else,for,while,class,
             import,export,default,new,this,static,extends,typeof,
             null,undefined,true,false,async,await},
  keywordstyle=\color{azulsecundario}\bfseries,
  ndkeywords={Math,Number,Array,Object,window,document,console},
  ndkeywordstyle=\color{teal},
  identifierstyle=\color{black},
  sensitive=true,
  comment=[l]{//},
  morecomment=[s]{/*}{*/},
  commentstyle=\color{gray}\itshape,
  stringstyle=\color{orange!80!black},
  morestring=[b]",
  morestring=[b]',
  morestring=[b]`,
}
\lstset{
  language=JavaScript,
  backgroundcolor=\color{griscode},
  basicstyle=\ttfamily\footnotesize,
  numbers=left,
  numberstyle=\tiny\color{gray},
  stepnumber=1,
  numbersep=8pt,
  frame=single,
  rulecolor=\color{gray!30},
  breaklines=true,
  breakatwhitespace=true,
  captionpos=b,
  autogobble=true,
  xleftmargin=15pt,
}

% ============================================================
%  ENCABEZADO Y PIE
% ============================================================
\pagestyle{fancy}
\fancyhf{}
\fancyhead[L]{\small\textcolor{gray}{ElectroLab-JS --- Simulador de Electromagnetismo}}
\fancyhead[R]{\small\textcolor{gray}{\nouppercase{\leftmark}}}
\fancyfoot[C]{\small\textcolor{gray}{\thepage}}
\renewcommand{\headrulewidth}{0.4pt}
\renewcommand{\footrulewidth}{0.4pt}

% ============================================================
%  MACROS ÚTILES
% ============================================================
\newcommand{\V}{\,\si{\volt}}
\newcommand{\A}{\,\si{\ampere}}
\newcommand{\W}{\,\si{\watt}}
\newcommand{\ohmunit}{\,\si{\ohm}}
\newcommand{\F}{\,\si{\farad}}
\newcommand{\T}{\,\si{\tesla}}
\newcommand{\Req}{R_{\mathrm{eq}}}
\newcommand{\Imax}{I_{\mathrm{max}}}
\newcommand{\Pmax}{P_{\mathrm{max}}}
\newcommand{\Vcrit}{V_{\mathrm{cr\acute{i}tico}}}
\newcommand{\Rcrit}{R_{\mathrm{cr\acute{i}tica}}}
\newcommand{\js}{\texttt{JavaScript}}
\newcommand{\mvc}{\textsc{mvc}}

% ============================================================
%  DOCUMENTO
% ============================================================
\begin{document}

% ?? Portada ??????????????????????????????????????????????????
\begin{titlepage}
  \centering
  \vspace*{1.5cm}

  {\Large\bfseries\color{azulprimario} UNIVERSIDAD TECNOLÓGICA NACIONAL}\\[0.3cm]
  {\large\color{gray} Facultad de Ingeniería en Sistemas}\\[0.15cm]
  {\normalsize\color{gray} Materia: Física para Ingeniería II}

  \vspace{1.5cm}
  \textcolor{azulsecundario}{\rule{\linewidth}{2pt}}\\[0.4cm]

  {\Huge\bfseries\color{azulprimario} ElectroLab-JS}\\[0.4cm]
  {\LARGE\bfseries\color{azulsecundario}
    Simulador Interactivo de Electricidad\\[0.2cm] y Magnetismo}\\[0.4cm]

  \textcolor{azulsecundario}{\rule{\linewidth}{2pt}}

  \vspace{1.2cm}

  \begin{tcolorbox}[
    colback=grisclaro, colframe=azulsecundario,
    width=0.7\linewidth, halign=center
  ]
    \centering
    \small
    \textbf{Alumno:} [Tu Nombre Completo]\\
    \textbf{Matrícula:} [Número de Matrícula]\\[6pt]
    \textbf{Docente:} Dr. [Nombre del Profesor]\\[6pt]
    \textbf{Proyecto Final de Curso}\\
    Abril, 2026
  \end{tcolorbox}

  \vfill
  {\footnotesize\color{gray}
    Tecnologías: HTML5 ? CSS3 ? JavaScript Vanilla ? Canvas API\\
    Arquitectura: MVC ? ES6 Modules ? Principios SOLID
  }
\end{titlepage}

% ?? Índice ???????????????????????????????????????????????????
\tableofcontents
\newpage

% ?? Resumen ??????????????????????????????????????????????????
\begin{abstract}
\noindent
El presente trabajo documenta el diseño, desarrollo e implementación de
\textbf{ElectroLab-JS}, una aplicación web interactiva orientada a la
enseñanza visual de conceptos fundamentales de electricidad y magnetismo a
nivel universitario. El simulador cubre tres módulos: la Ley de Ohm con
análisis de riesgo térmico, topologías de circuitos eléctricos (serie,
paralelo y mixto) con simulación de fallas, y el electromagnetismo basado
en la Fuerza de Lorentz aplicado a un motor de corriente continua. La
herramienta fue construida con tecnologías web nativas ?HTML5, CSS3,
\js{} Vanilla y Canvas API? sin dependencias externas, bajo el patrón
\mvc{} y principios SOLID. Los conceptos simulados tienen sustento directo
en las prácticas de laboratorio realizadas durante el semestre, incluyendo
electrización por frotamiento, verificación experimental de la Ley de
Coulomb y fabricación de capacitores artesanales. Los resultados
demuestran que la representación visual animada de fenómenos físicos
facilita la comprensión de conceptos abstractos como la resistencia
equivalente, la distribución de corriente y la inducción electromagnética.

\bigskip
\noindent\textbf{Palabras clave:} Ley de Ohm, circuitos eléctricos,
fuerza de Lorentz, capacitancia, electrización, simulación web, Canvas API,
\js{}, educación en ingeniería.
\end{abstract}

% ============================================================
\section{Introducción}
% ============================================================

\subsection{Planteamiento del problema}

La física eléctrica es una de las materias con mayor índice de dificultad
en carreras de ingeniería. La brecha entre la formulación matemática
abstracta ---ecuaciones, vectores, campos--- y la intuición que el
estudiante necesita para resolver problemas reales es una dificultad
ampliamente documentada en la literatura pedagógica~\cite{serway2018,
hayt2019}.

Ver una fórmula escrita en el pizarrón es fundamentalmente diferente a
\emph{ver} cómo los electrones se detienen completamente cuando un foco se
funde en un circuito en serie, o cómo un cable se dobla hacia un imán
cuando la corriente fluye. Las herramientas de simulación interactiva se
presentan como una alternativa pedagógica accesible y eficaz ante las
limitaciones del laboratorio físico: tiempo, equipo costoso y
reproducibilidad.

\subsection{Justificación}

La decisión de construir un simulador propio responde a tres razones
concretas. Primero, la oportunidad de implementar directamente las fórmulas
estudiadas en clase, verificando su comprensión real. Segundo, la
posibilidad de diseñar animaciones y metáforas visuales en torno a los
ejemplos vistos en el semestre. Tercero, el valor adicional de aprender
simultáneamente arquitectura de software, patrones de diseño y física
computacional.

\subsection{Objetivos}

\textbf{General:} Diseñar e implementar una aplicación web interactiva que
permita simular y visualizar en tiempo real los principios fundamentales de
la Ley de Ohm, las topologías de circuitos y el electromagnetismo.

\textbf{Específicos:}
\begin{itemize}[leftmargin=1.4em]
  \item Modelar la Ley de Ohm e implementar análisis de riesgo térmico.
  \item Calcular $\Req$ para circuitos serie, paralelo y mixto, con
        simulación de fallas.
  \item Ilustrar la Fuerza de Lorentz con física de simulación (Verlet) y
        su aplicación en un motor DC.
  \item Integrar los conceptos de las prácticas del semestre como marco
        conceptual del simulador.
  \item Aplicar el patrón \mvc{} y principios SOLID.
\end{itemize}

% ============================================================
\section{Marco teórico}
% ============================================================

\subsection{Electrización y carga eléctrica}

\subsubsection{Historia y contexto}

El estudio sistemático de la electricidad comenzó con William Gilbert en
1600. La distinción entre dos tipos de carga fue establecida por Charles
Fran?ois du Fay en 1733, y fue Benjamin Franklin quien estableció la
convención de carga positiva y negativa aún vigente~\cite{griffiths2017}.

\begin{cajaclas}
En el laboratorio se realizó la práctica de electrización por frotamiento
con materiales como acrílico, PVC, vidrio y cobre, frotados con paños de
poliéster, franela, peluche y hule espuma. Se registró atracción o
repulsión según el tipo de carga transferida, usando un sensor de rotación
Vernier conectado a un LabQuest para cuantificar la interacción.
\end{cajaclas}

\figplaceholder{Práctica 1 --- Electrización por frotamiento}
{Inserta aquí la fotografía del montaje experimental con los materiales y el sensor Vernier.}

\subsubsection{Ley de Coulomb}

En 1785, Charles-Augustin de Coulomb cuantificó la fuerza entre cargas
eléctricas puntuales:

\begin{equation}
  F = k \frac{q_1\, q_2}{r^2}, \qquad
  k = \SI{8.99e9}{\newton\metre\squared\per\coulomb\squared}
  \label{eq:coulomb}
\end{equation}

La relación inversa al cuadrado es compartida con la gravitación
newtoniana, reflejando una simetría profunda en las fuerzas de largo
alcance de la naturaleza.

\begin{cajaclas}
Se verificó experimentalmente la Ley de Coulomb midiendo la fuerza de
repulsión entre dos barras de acrílico cargadas a distintas distancias. El
ajuste logarítmico de la gráfica $F$ vs. $r$ arrojó una pendiente $\approx
-2$, confirmando $F \propto 1/r^2$ conforme predice la
ecuación~\eqref{eq:coulomb}.
\end{cajaclas}

\figplaceholder{Práctica 1 --- Gráfica F vs. r (Ley de Coulomb)}
{Inserta aquí la gráfica obtenida del LabQuest con el ajuste logarítmico.}

\subsection{Capacitores y almacenamiento de energía eléctrica}

La capacitancia $C$ cuantifica la capacidad de almacenamiento de carga de
un dispositivo a un dado voltaje:

\begin{equation}
  C = \frac{Q}{V}
  \label{eq:cap_def}
\end{equation}

Para un capacitor de placas paralelas con dieléctrico:

\begin{equation}
  C = \frac{\kappa \,\varepsilon_0\, A}{d}, \qquad
  \varepsilon_0 = \SI{8.85e-12}{\farad\per\metre}
  \label{eq:cap_placas}
\end{equation}

donde $\kappa$ es la constante dieléctrica, $A$ el área de las placas y
$d$ su separación. La energía almacenada es:

\begin{equation}
  E = \frac{1}{2}CV^2
  \label{eq:energia_cap}
\end{equation}

\begin{cajaclas}
Se fabricaron capacitores electrolíticos artesanales con placas de aluminio
y soluciones de borax y bicarbonato de sodio. Al aplicar \SI{12}{\V} al
ánodo de aluminio, se formó una capa de Al$_2$O$_3$ nanométrica que actúa
como dieléctrico. La capacitancia obtenida con borax fue
$C \approx \SI{10}{\micro\farad}$.
\end{cajaclas}

\figplaceholder{Práctica 2 --- Capacitor electrolítico artesanal}
{Inserta aquí la foto del vaso de precipitados con las placas de aluminio, la fuente y el multímetro en modo capacímetro.}

\subsection{Ley de Ohm y potencia eléctrica}

\subsubsection{Georg Simon Ohm (1789--1854)}

Georg Simon Ohm era maestro de secundaria cuando publicó en 1827 \emph{Die
galvanische Kette, mathematisch bearbeitet}, donde estableció la ley que
lleva su nombre~\cite{ohm1827}. Su trabajo fue inicialmente ridiculizado en
Alemania; no fue reconocido hasta que la Royal Society de Londres le
otorgó la Medalla Copley en 1841.

\begin{cajahist}
\textbf{Dato histórico:} Ohm usó termopares en lugar de pilas químicas
porque ofrecían un voltaje más estable. Esta precisión experimental fue
clave para obtener resultados reproducibles.
\end{cajahist}

\subsubsection{La ley}

Para un conductor óhmico a temperatura constante:

\begin{equation}
  I = \frac{V}{R}
  \label{eq:ohm}
\end{equation}

La resistencia de un conductor depende de sus propiedades geométricas y
del material:

\begin{equation}
  R = \rho \frac{L}{A}
  \label{eq:resistividad}
\end{equation}

\subsubsection{Potencia y Efecto Joule}

James Prescott Joule (1818--1889) estableció en 1841 la equivalencia entre
corriente eléctrica y calor. La potencia disipada adopta tres formas
equivalentes:

\begin{equation}
  P = VI = I^2 R = \frac{V^2}{R}
  \label{eq:potencia}
\end{equation}

De \eqref{eq:potencia} se extraen los umbrales críticos de operación segura:

\begin{align}
  \Vcrit &= \sqrt{P_{\max} \cdot R}  \label{eq:vcrit} \\
  \Rcrit &= \frac{V^2}{P_{\max}}     \label{eq:rcrit}
\end{align}

Estos valores se muestran en tiempo real en el panel de predicción del
simulador.

\begin{cajaclas}
El ejemplo del profesor con distintos valores de resistencia conectados a
una fuente variable demostró que duplicar el voltaje \emph{cuadruplica} la
potencia disipada (ecuación~\eqref{eq:potencia}). Esta relación cuadrática,
contraintuitiva para el estudiante, motivó que el slider de voltaje del
simulador esté directamente vinculado a la velocidad de los electrones y al
brillo de la resistencia de manera no lineal.
\end{cajaclas}

\subsection{Circuitos eléctricos}

\subsubsection{Kirchhoff y la teoría de circuitos}

Gustav Kirchhoff (1824--1887) formuló en 1845, con 21 años, las dos leyes
que sustentan el análisis de cualquier circuito de corriente continua:

\begin{itemize}
  \item \textbf{Primera ley (Nodos):} $\sum I_{\mathrm{entrada}} = \sum I_{\mathrm{salida}}$ \quad (conservación de carga)
  \item \textbf{Segunda ley (Mallas):} $\sum V_k = 0$ \quad (conservación de energía)
\end{itemize}

\subsubsection{Circuito en serie}

\begin{equation}
  \Req^{\mathrm{serie}} = \sum_{i=1}^{n} R_i, \qquad
  I_1 = I_2 = \cdots = I_n = I
  \label{eq:serie}
\end{equation}

La interrupción de cualquier nodo corta el flujo en toda la red:
$R_{\mathrm{falla}} \to \infty \Rightarrow \Req \to \infty \Rightarrow I = 0$.

\subsubsection{Circuito en paralelo}

\begin{equation}
  \frac{1}{\Req^{\mathrm{par}}} = \sum_{i=1}^{n} \frac{1}{R_i}, \qquad
  \Req < R_{\min}
  \label{eq:paralelo}
\end{equation}

Cada rama recibe el voltaje completo; la corriente total aumenta con cada
rama adicional.

\subsubsection{Circuito mixto --- algoritmo implementado}

El motor matemático recorre la secuencia de componentes de derecha a
izquierda acumulando $\Req$:

\begin{lstlisting}[caption={Nucleo del algoritmo de Req mixto (ModeloCircuitos.js)},
                   label={lst:req_mixto}]
for (let i = secuencia.length - 1; i >= 0; i--) {
    const tipo = secuencia[i];
    const r    = resEfectivas[i];

    if (tipo === 'S_top' || tipo === 'S_bot') {
        r_acc += r;                          // Serie: suma directa
    } else if (tipo === 'P') {
        if (r_acc === Infinity) r_acc = r;
        else r_acc = (r_acc * r) / (r_acc + r); // Paralelo: combinacion
    }
}
\end{lstlisting}

El tratamiento del foco fundido es elegante: asignar $R = \infty$ al
componente fallido hace que \texttt{JavaScript} propague correctamente
$\infty + k = \infty$ (bloquea serie) y $1/\infty = 0$ (elimina la rama
paralela).

\begin{cajaclas}
El ejemplo del profesor con una batería, tres focos en serie y un
interruptor que simulaba el fundido de un filamento fue el momento que
motivó que la simulación de fallas sea el elemento central del módulo de
circuitos. Al ver que apagar uno apagaba todos, y que en paralelo solo se
apagaba la habitación afectada, el concepto quedó grabado visualmente.
\end{cajaclas}

\subsection{Electromagnetismo y Fuerza de Lorentz}

\subsubsection{Historia: de Oersted a Maxwell}

El 21 de abril de 1820, Hans Christian Oersted descubrió durante una clase
que una corriente eléctrica desvía una aguja de brújula, naciendo así el
electromagnetismo. Ampère formalizó la relación matemáticamente; Faraday
descubrió la inducción en 1831; Maxwell la unificó todo en cuatro
ecuaciones diferenciales entre 1861 y 1865.

\begin{cajahist}
Hendrik Lorentz (1853--1928) fue el primer físico en recibir el Premio
Nobel de Física (1902, compartido con Pieter Zeeman). Su expresión de la
fuerza sobre una carga en movimiento unificó los efectos eléctrico y
magnético en una sola ecuación vectorial.
\end{cajahist}

\subsubsection{Fuerza de Lorentz}

Para una carga $q$ con velocidad $\mathbf{v}$ en campos $\mathbf{E}$ y
$\mathbf{B}$:

\begin{equation}
  \mathbf{F} = q(\mathbf{E} + \mathbf{v} \times \mathbf{B})
  \label{eq:lorentz_completa}
\end{equation}

Para un conductor de longitud efectiva $L$ con corriente $I$:

\begin{equation}
  \mathbf{F} = I\,\mathbf{L} \times \mathbf{B}, \qquad
  |\mathbf{F}| = I\,L\,B\,\sin\theta
  \label{eq:lorentz_cable}
\end{equation}

\subsubsection{Campo magnético generado por corriente (Ampère)}

\begin{equation}
  B = \frac{\mu_0\, I}{2\pi r}, \qquad
  \mu_0 = 4\pi \times 10^{-7} \;\si{\tesla\metre\per\ampere}
  \label{eq:ampere}
\end{equation}

\begin{cajaclas}
El experimento del profesor fue el detonador del módulo de magnetismo: al
pasar corriente por un cable flexible y acercar un imán de neodimio, el
cable se desplazaba visiblemente. Al invertir la corriente o girar el imán,
la fuerza se invertía. Sin corriente, el cable no respondía. El módulo de
magnetismo del simulador reproduce exactamente este experimento con
interacción de arrastre del imán.
\end{cajaclas}

\figplaceholder{Demostración de la Fuerza de Lorentz}
{Inserta aquí la foto del experimento del profesor: cable flexible, fuente de corriente e imán de neodimio.}

\subsubsection{Integración de Verlet para la simulación del cable}

El cable se modela como una cadena de $N = 20$ nodos con restricciones de
distancia. La posición de cada nodo libre se actualiza con el integrador de
Verlet~\cite{verlet1967}:

\begin{equation}
  x_{n+1} = 2x_n - x_{n-1} + a \cdot \Delta t^2
  \label{eq:verlet}
\end{equation}

donde $a = F_{\mathrm{Lorentz}} / m$ y un factor de amortiguación de $0.95$
por paso previene oscilaciones divergentes. Este método conserva mejor la
energía que la integración de Euler simple, produciendo movimientos
visualmente más naturales.

% ============================================================
\section{Casos de uso en física}
% ============================================================

Los siguientes casos de uso describen escenarios reales donde los
principios del semestre tienen aplicación directa y medible.

\subsection{Caso 1: Instalación eléctrica doméstica}

\begin{casodeuso}{Instalación eléctrica doméstica}
\textbf{Concepto:} Ley de Ohm + Circuitos en paralelo.\\[4pt]
\textbf{Descripción:} En una vivienda todos los circuitos se conectan en
paralelo respecto a la línea de \SI{127}{\V}. Cada dispositivo opera
independientemente, recibiendo el voltaje nominal completo.\\[4pt]
\textbf{Ejemplo numérico:}
\begin{itemize}
  \item Foco LED \SI{60}{\W}: $R = V^2/P = 127^2/60 = \SI{268.8}{\ohm}$, $I = \SI{0.47}{\A}$
  \item Televisor \SI{120}{\W}: $I = \SI{0.94}{\A}$
  \item Cafetera \SI{900}{\W}: $I = \SI{7.09}{\A}$
  \item $I_{\mathrm{total}} = \SI{8.50}{\A}$ (seguro, $< \SI{15}{\A}$)
  \item Plancha adicional \SI{1500}{\W}: $I_{\mathrm{plancha}} = \SI{11.81}{\A}$ $\Rightarrow$ $I_{\mathrm{total}} = \SI{20.31}{\A} > \SI{15}{\A}$ $\Rightarrow$ \textbf{breaker dispara}
\end{itemize}
\textbf{Importancia:} Dimensionamiento correcto de cables y protecciones para prevenir incendios.
\end{casodeuso}

\vspace{4pt}
\begin{table}[H]
\centering\small
\caption{Resumen --- Caso de uso 1}
\begin{tabularx}{0.85\linewidth}{>{\bfseries}lX}
\toprule
Campo & Contenido \\ \midrule
Nombre       & Instalación eléctrica doméstica \\
Concepto     & Ley de Ohm + Circuitos en paralelo \\
Contexto     & Viviendas, oficinas, instalaciones comerciales \\
Aplicación   & Dimensionamiento de cables y protecciones \\
Fórmula clave & $I_{\mathrm{total}} = V/\Req$,\;\; $1/\Req = \sum 1/R_i$ \\
Impacto      & Prevención de incendios por sobrecalentamiento de cables \\ \bottomrule
\end{tabularx}
\end{table}

\subsection{Caso 2: Carga de smartphones y Efecto Joule}

\begin{casodeuso}{Carga segura de smartphones}
\textbf{Concepto:} Ley de Ohm + Potencia cuadrática + Efecto Joule.\\[4pt]
\textbf{Descripción:} Un cargador genérico de alto voltaje puede destruir
el PMIC del teléfono porque la potencia disipada crece con el cuadrado del
voltaje.\\[4pt]
\textbf{Ejemplo numérico (iPhone 5, $V_{\max} = \SI{5}{\V}$):}
\begin{itemize}
  \item Cargador correcto (\SI{5}{\V}): $P_1 = 25/10 = \SI{2.5}{\W}$
  \item Cargador incorrecto (\SI{9}{\V}): $P_2 = 81/10 = \SI{8.1}{\W} = 3.24\times P_1$
\end{itemize}
\textbf{Importancia:} Base del protocolo USB-PD que negocia el voltaje
dinámicamente para proteger el dispositivo.
\end{casodeuso}

\vspace{4pt}
\begin{table}[H]
\centering\small
\caption{Resumen --- Caso de uso 2}
\begin{tabularx}{0.85\linewidth}{>{\bfseries}lX}
\toprule
Campo & Contenido \\ \midrule
Nombre       & Carga segura de smartphones \\
Concepto     & Efecto Joule, potencia eléctrica, Ley de Ohm \\
Contexto     & Electrónica de consumo, cargadores USB \\
Aplicación   & Compatibilidad de cargadores, protocolo USB-PD \\
Fórmula clave & $P = V^2/R$ \;\;(cuadrática en $V$) \\
Impacto      & Protección de dispositivos y diseño de estándares de carga \\ \bottomrule
\end{tabularx}
\end{table}

\subsection{Caso 3: Foco fundido --- serie vs. paralelo}

\begin{casodeuso}{Foco fundido: serie vs. paralelo (6 focos, $R = \SI{240}{\ohm}$, $V = \SI{120}{\V}$)}
\textbf{En serie:}
$\Req = 6 \times 240 = \SI{1440}{\ohm}$,\; $I = \SI{0.083}{\A}$.
Foco fundido: $\Req \to \infty$,\; $I = 0$. Toda la casa se apaga.\\[4pt]
\textbf{En paralelo:}
$\Req = 240/6 = \SI{40}{\ohm}$,\; $I_{\mathrm{total}} = \SI{3}{\A}$.
Foco fundido: $\Req = 240/5 = \SI{48}{\ohm}$,\; $I_{\mathrm{total}} = \SI{2.5}{\A}$.
Solo esa rama se apaga.\\[4pt]
\textbf{Importancia:} Explica por qué la instalación doméstica moderna usa
exclusivamente topología en paralelo, y por qué las luces LED navideñas no
se apagan todas al fundirse una.
\end{casodeuso}

\subsection{Caso 4: Electrización y Ley de Coulomb en materiales}

\begin{casodeuso}{Electrización y Ley de Coulomb (Práctica de laboratorio)}
\textbf{Concepto:} Carga eléctrica + Serie triboeléctrica + $F = kq_1q_2/r^2$.\\[4pt]
\textbf{Datos experimentales:}
\begin{itemize}
  \item Material: barra de acrílico frotada con poliéster.
  \item A $r_2 = 2r_1$: $F_2 \approx F_1/4$ (verificado con sensor Vernier).
  \item Ajuste logarítmico: pendiente $\approx -2$, confirmando $F \propto 1/r^2$.
\end{itemize}
\textbf{Impacto:} Base de precipitadores electrostáticos, pintado
electrostático e impresoras láser.
\end{casodeuso}

\subsection{Caso 5: Capacitor electrolítico --- flash de cámara}

\begin{casodeuso}{Almacenamiento de energía en capacitor (Práctica de laboratorio)}
\textbf{Concepto:} $C = \kappa\varepsilon_0 A/d$,\; $E = \tfrac{1}{2}CV^2$.\\[4pt]
\textbf{Flash de cámara:} $C = \SI{1000}{\micro\farad}$, $V = \SI{300}{\V}$
$\Rightarrow E = \SI{45}{\joule}$.
Descarga en $\SI{1}{\milli\second}$ $\Rightarrow P_{\mathrm{pico}} = \SI{45}{\kilo\watt}$.\\[4pt]
\textbf{Capacitor artesanal (borax):} $C \approx \SI{10}{\micro\farad}$, $V = \SI{12}{\V}$
$\Rightarrow E = \SI{720}{\micro\joule}$.\\[4pt]
\textbf{Impacto:} La práctica conecta directamente la ecuación
\eqref{eq:cap_placas} con la experiencia de ver crecer la capa de óxido y
medir $C$ con el multímetro.
\end{casodeuso}

\subsection{Caso 6: Motor DC en sistemas de bombeo industrial}

\begin{casodeuso}{Motor DC en bomba de agua}
\textbf{Concepto:} $F = IL \times B$,\; $\omega \propto K_v V$,\; calentamiento por Joule.\\[4pt]
\textbf{Parámetros del motor simulado:}
$K_v = \SI{60}{RPM\per\volt}$,\; $R_{\mathrm{int}} = \SI{1.5}{\ohm}$,\;
$T_{\mathrm{cr}} = \SI{130}{\celsius}$.\\[4pt]
\textbf{Sobrevoltaje a \SI{55}{\V}:} exceso $= 55 - 48 = \SI{7}{\V}$;
calentamiento $\propto 7^2 = 49$ veces más rápido que en operación
marginal. Tiempo estimado hasta falla: $< \SI{30}{\second}$.\\[4pt]
\textbf{Impacto:} Base del diseño de protecciones térmicas y controladores
PWM en industria y electromovilidad.
\end{casodeuso}

\subsection{Caso 7: Cable en campo magnético (experimento de clase)}

\begin{casodeuso}{Cable en campo magnético externo}
\textbf{Concepto:} $F = I\,L\,B\,\sin\theta$,\; Regla de la Mano Derecha.\\[4pt]
\textbf{Estimación experimental:}
$I = \SI{5}{\A}$,\; $L = \SI{0.10}{\metre}$,\; $B \approx \SI{0.05}{\tesla}$ (imán neodimio a \SI{2}{\centi\metre})
$\Rightarrow F = \SI{25}{\milli\newton}$ (suficiente para deflexión visible).\\[4pt]
Al duplicar $I \to \SI{10}{\A}$: $F \to \SI{50}{\milli\newton}$, desplazamiento duplicado.\\[4pt]
\textbf{Impacto:} Fundamento de altavoces, galvanómetros, actuadores
electromagnéticos y todos los motores eléctricos.
\end{casodeuso}

% ============================================================
\section{Proceso de diseño y desarrollo}
% ============================================================

\subsection{Decisiones tecnológicas}

\subsubsection{JavaScript Vanilla sin frameworks}

La elección de \js{} Vanilla responde a tres criterios:

\begin{enumerate}
  \item \textbf{Accesibilidad universal:} funciona en cualquier navegador
        sin instalación ni proceso de compilación.
  \item \textbf{Transparencia matemática:} sin capas de abstracción, la
        línea \texttt{this.resultados.req = resEfectivas.reduce((a,b) =>
        a+b, 0)} \emph{es} literalmente el cálculo de $\Req$ en serie.
  \item \textbf{Rendimiento nativo:} redibujo de 80 partículas a 60 fps
        sin overhead de estado reactivo.
\end{enumerate}

\subsubsection{Canvas API vs. SVG y WebGL}

Canvas API fue elegido sobre SVG por rendimiento: SVG mantiene un árbol
DOM que se degrada con muchos elementos animados. WebGL ofrece mayor
rendimiento pero con complejidad desproporcionada al objetivo pedagógico.

\subsubsection{Integración de Verlet vs. Euler}

El método de Euler acumula error en cada paso, produciendo oscilaciones
crecientes en sistemas con fuerzas restauradoras. Verlet deriva la
velocidad implícitamente de la diferencia de posiciones
(ecuación~\eqref{eq:verlet}), conservando mejor la energía con el mismo
costo computacional~\cite{verlet1967}.

\subsection{Arquitectura MVC}

\begin{table}[H]
\centering\small
\caption{Distribución de responsabilidades por capa MVC}
\begin{tabularx}{\linewidth}{>{\bfseries}llX}
\toprule
Capa & Clase principal & Responsabilidad \\ \midrule
Modelo      & \texttt{ModeloOhm}, \texttt{ModeloCircuitos}  & Estado, cálculos físicos, notificaciones \\
Vista       & \texttt{VistaOhm}, \texttt{VistaCircuitos}    & HTML dinámico, referencias DOM, Canvas \\
Controlador & \texttt{ControladorOhm}, \texttt{ControladorCircuitos} & Eventos, bucle animación, ciclo de vida \\
Compartido  & \texttt{RenderizadorBase}, \texttt{AnimadorElectrones} & Componentes reutilizables entre módulos \\
\bottomrule
\end{tabularx}
\end{table}

La comunicación Modelo $\leftrightarrow$ Vista ocurre mediante un observador
simple: \texttt{modelo.alActualizar = () => \{ vista.actualizarUI(modelo) \}}.
El bucle \texttt{requestAnimationFrame} vive en el Controlador y se cancela al
cambiar de módulo, evitando fugas de memoria.

\subsection{Bocetos y proceso iterativo}

\figplaceholder{Mockup 1 --- Módulo de Ley de Ohm}
{Inserta aquí el boceto inicial con la disposición del canvas, controles y panel de predicción.}

\figplaceholder{Mockup 2 --- Módulo de Circuitos}
{Inserta aquí el boceto del diagrama técnico y la transición al mundo real.}

\figplaceholder{Mockup 3 --- Módulo de Magnetismo}
{Inserta aquí el boceto del laboratorio de Lorentz y del motor DC.}

% ============================================================
\section{Resultados --- el simulador}
% ============================================================

\subsection{Módulo 1: Ley de Ohm}

El módulo implementa un circuito con una fuente de voltaje y una
resistencia, con electrones animados a velocidad proporcional a $I =
V/R$. El panel de predicción muestra $\Vcrit$ y $\Rcrit$ en tiempo real
(ecuaciones~\eqref{eq:vcrit} y \eqref{eq:rcrit}).

\figplaceholder{Captura 1 --- Módulo Ley de Ohm, estado normal}
{Electrones cian fluyendo, resistencia beige, panel de predicción visible.}

\figplaceholder{Captura 2 --- Estado crítico y quemado}
{Cable naranja brillante / resistencia rota con chispas.}

\figplaceholder{Captura 3 --- Modo carga de celulares}
{Teléfono cargando y teléfono dañado por exceso de voltaje.}

\subsection{Módulo 2: Circuitos}

\figplaceholder{Captura 4 --- Circuito en paralelo (diagrama técnico)}
{5 ramas en paralelo con electrones distribuyéndose y valores de Req e I.}

\figplaceholder{Captura 5 --- Mundo Real, foco fundido en serie}
{Todos los focos apagados; solo el fundido muestra la X roja.}

\figplaceholder{Captura 6 --- Mundo Real, foco fundido en paralelo}
{Solo una habitación apagada; el resto iluminado con normalidad.}

\subsection{Módulo 3: Magnetismo}

\figplaceholder{Captura 7 --- Laboratorio de Lorentz}
{Cable deflectado, líneas de campo animadas y vectores de fuerza verdes.}

\figplaceholder{Captura 8 --- Motor DC, operación normal y sobrevoltaje}
{Agua cayendo de la regadera vs. motor rojo con temperatura crítica.}

\figplaceholder{Captura 9 --- Vista zoom interior del motor}
{Rotor con bobinas brillando, imanes del estátor y conmutador central.}

\subsection{Verificación matemática}

\begin{table}[H]
\centering\small
\caption{Verificación de resultados del simulador contra valores teóricos}
\begin{tabular}{llll}
\toprule
Caso de prueba & {Valor teórico} & {Valor simulado} & Error \\
\midrule
$V=12\V,\,R=\SI{100}{\ohm} \to I$           & 0.12\,A  & 0.12\,A  & 0.00\% \\
$V=12\V,\,R=\SI{100}{\ohm} \to P$           & 1.44\,W  & 1.44\,W  & 0.00\% \\
$\Vcrit$ ($R=\SI{100}{\ohm},\,\Pmax=50\W$)  & 70.71\,V & 70.71\,V & 0.00\% \\
$\Req$: 3?\SI{240}{\ohm}\ en serie          & \SI{720}{\ohm} & \SI{720}{\ohm} & 0.00\% \\
$\Req$: 3?\SI{240}{\ohm}\ en paralelo       & \SI{80}{\ohm}  & \SI{80}{\ohm}  & 0.00\% \\
$I$ a 120\V\ con $\Req=\SI{80}{\ohm}$       & 1.50\,A  & 1.50\,A  & 0.00\% \\
Foco fundido en serie ($R_{eq}\to\infty$) & $I=0\,\si{\ampere}$ & $I=0\,\si{\ampere}$ & Correcto \\
Foco fundido en paralelo (5/6) & $I{-}17\,\%$ & $I{-}17\,\%$ & Correcto \\
\bottomrule
\end{tabular}
\end{table}

% ============================================================
\section{Conclusiones}
% ============================================================

El desarrollo de ElectroLab-JS fue simultáneamente un ejercicio de
comprensión física y de ingeniería de software. Implementar las fórmulas
de la Ley de Ohm, la resistencia equivalente de circuitos mixtos y la
física de Verlet para el cable electromagnético requirió entenderlas a un
nivel que va más allá de memorizarlas: fue necesario razonar sobre casos
límite ($R = 0$, $R = \infty$), sobre el comportamiento numérico de las
ecuaciones y sobre cómo traducir conceptos continuos en algoritmos
discretos que se ejecutan 60 veces por segundo.

La conexión con las prácticas del semestre resultó el elemento más valioso
del diseño pedagógico. Haber observado físicamente la deflexión del cable
ante el imán hizo que diseñar la simulación de la Fuerza de Lorentz fuera
una traducción directa de la experiencia al código. De igual manera,
fabricar un capacitor artesanal hizo tangibles los parámetros de la
ecuación~\eqref{eq:cap_placas}.

El valor educativo del simulador tiene una dimensión no prevista: fundir
un foco en serie y ver que toda la casa se apaga entrega una consecuencia
que en la vida real tomaría horas de trabajo eléctrico reproducir.

Como trabajo futuro se identifican tres extensiones de alto valor
pedagógico: circuitos de corriente alterna con representación fasorial,
un módulo de electrostática con visualización de campo eléctrico, y la
incorporación del módulo de capacitores como extensión natural del
laboratorio del semestre.

% ============================================================
%  REFERENCIAS
% ============================================================
\begin{thebibliography}{99}

\bibitem{serway2018}
R.~A. Serway y J.~W. Jewett,
\textit{Physics for Scientists and Engineers}, 10th~ed.
Cengage Learning, 2018.

\bibitem{hayt2019}
W.~H. Hayt, J.~E. Kemmerly y S.~M. Durbin,
\textit{Engineering Circuit Analysis}, 9th~ed.
McGraw-Hill Education, 2019.

\bibitem{nilsson2020}
J.~W. Nilsson y S.~A. Riedel,
\textit{Electric Circuits}, 11th~ed.
Pearson, 2020.

\bibitem{griffiths2017}
D.~J. Griffiths,
\textit{Introduction to Electrodynamics}, 4th~ed.
Cambridge University Press, 2017.

\bibitem{martin2017}
R.~C. Martin,
\textit{Clean Architecture: A Craftsman's Guide to Software Structure and Design}.
Prentice Hall, 2017.

\bibitem{gamma1994}
E.~Gamma, R.~Helm, R.~Johnson y J.~Vlissides,
\textit{Design Patterns: Elements of Reusable Object-Oriented Software}.
Addison-Wesley, 1994.

\bibitem{verlet1967}
L.~Verlet,
``Computer experiments on classical fluids. I. Thermodynamical properties of Lennard-Jones molecules,''
\textit{Physical Review}, vol.~159, no.~1, pp.~98--103, 1967.

\bibitem{mdn_canvas}
MDN Web Docs,
``Canvas API,''
Mozilla Developer Network, 2024.
[En línea]. Disponible: \url{https://developer.mozilla.org/en-US/docs/Web/API/Canvas_API}

\bibitem{usb_pd}
USB Implementers Forum,
\textit{USB Power Delivery Specification Revision 3.1}.
USB-IF, 2021.

\bibitem{nom001}
Secretaría de Energía,
\textit{Norma Oficial Mexicana NOM-001-SEDE-2012, Instalaciones Eléctricas (utilización)}.
México, 2012.

\bibitem{ohm1827}
G.~S. Ohm,
\textit{Die galvanische Kette, mathematisch bearbeitet}.
T.~H. Riemann, Berlín, 1827.

\bibitem{iie2019}
Instituto de Investigaciones Eléctricas,
\textit{Guía Técnica de Instalaciones Eléctricas Residenciales}.
IIE, Cuernavaca, México, 2019.

\end{thebibliography}

\end{document}
